% Options for packages loaded elsewhere
\PassOptionsToPackage{unicode}{hyperref}
\PassOptionsToPackage{hyphens}{url}
\documentclass[
  11pt,
]{article}
\usepackage{xcolor}
\usepackage[margin=2.4cm]{geometry}
\usepackage{amsmath,amssymb}
\setcounter{secnumdepth}{5}
\usepackage{iftex}
\ifPDFTeX
  \usepackage[T1]{fontenc}
  \usepackage[utf8]{inputenc}
  \usepackage{textcomp} % provide euro and other symbols
\else % if luatex or xetex
  \usepackage{unicode-math} % this also loads fontspec
  \defaultfontfeatures{Scale=MatchLowercase}
  \defaultfontfeatures[\rmfamily]{Ligatures=TeX,Scale=1}
\fi
\usepackage{lmodern}
\ifPDFTeX\else
  % xetex/luatex font selection
\fi
% Use upquote if available, for straight quotes in verbatim environments
\IfFileExists{upquote.sty}{\usepackage{upquote}}{}
\IfFileExists{microtype.sty}{% use microtype if available
  \usepackage[]{microtype}
  \UseMicrotypeSet[protrusion]{basicmath} % disable protrusion for tt fonts
}{}
\makeatletter
\@ifundefined{KOMAClassName}{% if non-KOMA class
  \IfFileExists{parskip.sty}{%
    \usepackage{parskip}
  }{% else
    \setlength{\parindent}{0pt}
    \setlength{\parskip}{6pt plus 2pt minus 1pt}}
}{% if KOMA class
  \KOMAoptions{parskip=half}}
\makeatother
\usepackage{longtable,booktabs,array}
\newcounter{none} % for unnumbered tables
\usepackage{calc} % for calculating minipage widths
% Correct order of tables after \paragraph or \subparagraph
\usepackage{etoolbox}
\makeatletter
\patchcmd\longtable{\par}{\if@noskipsec\mbox{}\fi\par}{}{}
\makeatother
% Allow footnotes in longtable head/foot
\IfFileExists{footnotehyper.sty}{\usepackage{footnotehyper}}{\usepackage{footnote}}
\makesavenoteenv{longtable}
\setlength{\emergencystretch}{3em} % prevent overfull lines
\providecommand{\tightlist}{%
  \setlength{\itemsep}{0pt}\setlength{\parskip}{0pt}}
\usepackage{amsmath,amssymb,amsthm}
\newtheorem{theorem}{Theorem}[section]
\newtheorem{lemma}[theorem]{Lemma}
\newtheorem{proposition}[theorem]{Proposition}
\newtheorem{corollary}[theorem]{Corollary}
\newtheorem{conjecture}[theorem]{Conjecture}
\theoremstyle{definition}
\newtheorem{definition}[theorem]{Definition}
\theoremstyle{remark}
\newtheorem{remark}[theorem]{Remark}
\usepackage{bookmark}
\IfFileExists{xurl.sty}{\usepackage{xurl}}{} % add URL line breaks if available
\urlstyle{same}
\hypersetup{
  pdftitle={Uniformly most split reliable graphs: a complete answer outside two edge counts},
  pdfauthor={Anonymous},
  hidelinks,
  pdfcreator={LaTeX via pandoc}}

\title{Uniformly most split reliable graphs: a complete answer outside
two edge counts\thanks{Research package:
https://doi.org/10.5281/zenodo.23450265}}
\author{Anonymous}
\date{18 October 2026 · Version 1.0.0 · Unrefereed candidate}

\begin{document}
\maketitle

\section*{Abstract}\label{abstract}
\addcontentsline{toc}{section}{Abstract}

A network with two distinguished terminals is split successfully when
independent edge failures leave exactly two connected components, one
containing each terminal. We ask when a connected simple graph with
fixed numbers of vertices and edges maximises this probability for every
common edge-retention probability. For \(n\ge7\), we prove nonexistence
for \(n\le m\le\binom n2-n-1\) and existence for \(m\ge\binom n2-n+2\);
trees admit an optimiser by earlier work. Only two edge counts per order
remain conditional. We reduce those cases to an explicit family of
cone-reliability inequalities and certify the required inequalities in
exact integer arithmetic through cone order \(50\), completing the
existence classification for \(7\le n\le51\). The structural argument
combines a prior classification of local optimisers with explicit
spanning-tree competitors and compatible extrema in the identity
\(\mathrm{SR}_{s,t}(H)=\mathrm{Rel}(H/st)-\mathrm{Rel}(H)\). We also
prove the boundary cone comparison at every order for \(p\ge1/2\) and on
an explicit positive interval near zero. The intervening probability
range remains open at unrestricted orders. We distinguish the universal
written proofs, the finite certificates needed for the boundary result,
and supplementary computational corroboration.

\section{Introduction}\label{introduction}

\subsection{Why this optimisation problem is
different}\label{why-this-optimisation-problem-is-different}

Ordinary network reliability rewards keeping every vertex connected.
Here the desired outcome is different: two distinguished sources must
remain separated, but every other vertex must still be connected to
exactly one of them. A surviving network with three components fails
this requirement, as does one connected component containing both
sources. The model is an idealised independent-edge failure model, not a
validated design prescription for a physical network.

The word \emph{uniformly} is the demanding part. We fix both the number
of vertices and the number of edges, and ask whether one design is best
for \textbf{every} common edge-retention probability. A design that wins
when most edges survive can lose when most fail. Consequently,
nonexistence of a uniform optimiser does not mean that no best design
exists at a specified probability; it means that the best choice cannot
be made independently of that probability.

Our proof separates these competing endpoint requirements. Near complete
edge failure, spanning two-forests determine the first nonzero term.
Near perfect edge retention, a prior classification forces a particular
split polynomial. In the intermediate density range these requirements
conflict. In the dense range, identifying the two sources converts the
comparison into two compatible all-terminal reliability comparisons.

\subsection{Problem and prior results}\label{problem-and-prior-results}

Let \(H\) be a connected simple graph with \(n\) vertices and \(m\)
edges and let \(s\ne t\) be two of its vertices, the terminals. Each
edge is deleted independently with probability \(1-p\). The
\emph{all-terminal reliability} \(\mathrm{Rel}(H;p)\) is the probability
that the surviving spanning subgraph is connected. The \emph{split
reliability} \(\mathrm{SR}_{s,t}(H;p)\), introduced by Brown and Mol
(2017) and studied as a uniform optimisation problem by Brown and
McMullin (2023), is the probability that the surviving spanning subgraph
has exactly two connected components, one containing \(s\) and the other
containing \(t\): every vertex can still communicate with exactly one of
the two terminals. Writing \(N_i=N_i(H)\) for the number of such
\emph{split states} with exactly \(i\) surviving edges,
\[\mathrm{SR}_{s,t}(H;p)=\sum_{i=n-2}^{m-1}N_i(H)\,p^i(1-p)^{m-i}.\] Let
\(\mathcal T_{n,m}\) be the class of two-terminal connected simple
graphs with \(n\) vertices and \(m\) edges. A member \(G\) is
\emph{uniformly most split reliable} (UMSR) in \(\mathcal T_{n,m}\) if
\(\mathrm{SR}_G(p)\ge\mathrm{SR}_H(p)\) for every
\(H\in\mathcal T_{n,m}\) and every \(p\in[0,1]\). Brown and McMullin
asked for which pairs \((n,m)\) such a graph exists. They showed that a
UMSR graph always exists for \(m=n-1\) (the path with terminals at its
ends), for \(m=\binom n2\) and for \(m=\binom n2-1\), that none exists
for \(m=n\ge 6\), and by computation that for \(n\le 5\) one exists for
every \(m\), for \(n=6\) one exists unless \(m\in\{6,8\}\), and for
\(n=7\) one exists if and only if \(m=6\) or \(14\le m\le 21\). The
remaining cases were left open. Romero (2025) characterised the
\emph{locally} most split reliable graphs (those optimal for \(p\) close
to \(1\)): they are exactly the two-terminal graphs with the same split
reliability polynomial as a two-terminal balloon graph \(G_{n,m}\).
Consequently any UMSR graph has the split polynomial of \(G_{n,m}\). He
proved that no UMSR graph exists when \(n\ge 7\) and
\(n\le m\le\binom{n-3}{2}+3\), and asked about the range
\(\binom{n-3}2+4\le m\le\binom n2-2\).

Our main result answers this for all but two values of \(m\) for each
\(n\), and for those two values reduces the question to a well known
conjecture in all-terminal reliability.

\begin{theorem}[Main theorem]\label{thm:main}
Let $n\ge 7$ and $n-1\le m\le\binom n2$, and write $k=\binom n2-m$ for the number of missing edges.
\begin{enumerate}
\item[(i)] If $m=n-1$ then a UMSR graph exists (the path with the terminals at its ends).
\item[(ii)] If $n\le m\le\binom n2-n-1$, equivalently $n+1\le k\le\binom n2-n$, then no UMSR graph exists in $\mathcal T_{n,m}$.
\item[(iii)] If $m\ge\binom n2-n+2$, equivalently $k\le n-2$, then the two-terminal balloon graph $G_{n,m}$ (Definition \ref{def:balloon}) is UMSR.
\item[(iv)] If $m\in\{\binom n2-n,\binom n2-n+1\}$, equivalently $k\in\{n-1,n\}$, then $G_{n,m}$ is UMSR provided the balloon graph $B_{n,m}$ is uniformly least reliable among connected graphs with $n$ vertices and $m$ edges (in the all-terminal sense). This proviso holds for $7\le n\le 51$, and for every $n\ge 11$ it follows from inequality $(\star)_{n-1,\varepsilon}$ of the boundary-reduction subsection (Theorem \ref{thm:reduction}).
\end{enumerate}
In particular, for $7\le n\le 51$ a UMSR graph exists in $\mathcal T_{n,m}$ if and only if $m=n-1$ or $m\ge\binom n2-n$.
\end{theorem}

For \(7\le n\le51\), the classification of vertex-and-edge counts is
simple: a uniform optimiser exists precisely for trees or when at most
\(n\) edges are missing. For larger orders the two boundary edge counts
remain conditional. This is an existence classification, not a claim
that every graph in an admissible class is optimal. Part (ii) is new for
\(\binom{n-3}2+4\le m\le\binom n2-n-1\) (for smaller \(m\) it is
Romero's Theorem 12). Part (iii) is new for \(k\ge 2\). The two edge
counts in (iv) are the two values immediately below the known dense
range for the all-terminal ``uniformly least reliable graph'' question:
Petingi, Saccoman and Schoppmann (1996) proved that balloon graphs are
uniformly least reliable for \(m\ge\binom{n-1}2+1\) (that is,
\(k\le n-2\)), and Romero (2022) proved the conjecture of Boesch,
Satyanarayana and Suffel for \(m\le n+3\); the values \(m=\binom{n-1}2\)
and \(\binom{n-1}2-1\) lie just below the dense range. Our reduction
shows that progress on that conjecture at these two edge counts
immediately completes the split reliability picture, and Theorem
\ref{thm:reduction} reduces the conjecture at these two edge counts, for
all \(n\), to a family of inequalities between explicit cone graphs,
which we verify exactly for \(n\le51\).

The reader can follow three proof stages. First, the preliminaries
establish the identification identity and exchange tool. Next, the
nonexistence argument uses Romero's candidate theorem and two explicit
competitors. Finally, the dense-existence argument combines
identified-graph maximality with an imported all-terminal minimum
theorem, then isolates the two boundary cases. The computation section
distinguishes the exact finite certificates from supplementary checks.
Readers interested chiefly in the scope can read the main theorem, the
small triangle example, and the final discussion before the technical
proofs.

\section{Preliminaries}\label{preliminaries}

The following notation separates three objects that must not be
conflated: the original two-terminal graph, its terminal identification,
and the ordinary balloon graph used in a least-reliability theorem.

{\def\LTcaptype{none} % do not increment counter
\begin{longtable}[]{@{}
  >{\raggedright\arraybackslash}p{(\linewidth - 4\tabcolsep) * \real{0.3333}}
  >{\raggedright\arraybackslash}p{(\linewidth - 4\tabcolsep) * \real{0.3333}}
  >{\raggedright\arraybackslash}p{(\linewidth - 4\tabcolsep) * \real{0.3333}}@{}}
\toprule\noalign{}
\begin{minipage}[b]{\linewidth}\raggedright
Symbol
\end{minipage} & \begin{minipage}[b]{\linewidth}\raggedright
Meaning
\end{minipage} & \begin{minipage}[b]{\linewidth}\raggedright
Role in the proof
\end{minipage} \\
\midrule\noalign{}
\endhead
\bottomrule\noalign{}
\endlastfoot
\(H,s,t\) & Connected simple graph and its two distinct terminals &
Competing network \\
\(n,m,k\) & Vertices, edges, and \(k=\binom n2-m\) missing edges & Fixed
comparison class \\
\(p,q\) & Edge retention and failure probabilities, \(q=1-p\) & Uniform
comparison variable \\
\(H/st\) & Terminal identification, with parallel edges retained & First
term of the identity \\
\(B_{n,m}\) & Ordinary balloon graph & All-terminal minimisation \\
\(G_{n,m}\) & Balloon with diameter terminals & Split-reliability
candidate \\
\(N_i,c_i\) & Split-state and connected-state edge counts & Exact finite
certificates \\
\(C_N(s)\) & Clique on \(N-1\) vertices plus a vertex with \(s\)
neighbours & Boundary reduction \\
\end{longtable}
}

\subsection{Notation}\label{notation}

Graphs are finite; simple unless stated otherwise. For a loopless
multigraph \(Y\) we write \(\mu_Y(uv)\) for the number of \(uv\) edges
and \(e_v(U)=\sum_{u\in U}\mu_Y(vu)\) for the number of edges (with
multiplicity) from \(v\) into a vertex set \(U\). \(K_q\) is the
complete graph, \(K_{1,a}\) the star with \(a\) leaves, \(K_2\) a single
edge, and \(K_q-F\) denotes \(K_q\) with the edge set of a graph \(F\)
on (a subset of) its vertices removed. \(\tau(Y)\) is the number of
spanning trees of \(Y\) and \(c_i(Y)\) the number of connected spanning
subgraphs with exactly \(i\) edges, so that
\(\mathrm{Rel}(Y;p)=\sum_i c_i(Y)p^i(1-p)^{|E(Y)|-i}\) and
\(\mathrm{Rel}(Y;p)=\tau(Y)p^{|V(Y)|-1}+O(p^{|V(Y)|})\) as \(p\to 0\).
The \emph{two-terminal} version of a graph carries an unordered pair
\(\{s,t\}\) of distinct terminals; the split polynomial is symmetric in
\(s\) and \(t\).

A \emph{split-\((s,t)\)-bridge} of \(H\) is an edge whose deletion
leaves \(s\) and \(t\) in different components. Since \(N_{m-1}(H)\)
counts single edges whose deletion produces a split state,
\(N_{m-1}(H)\) equals the number of split-\((s,t)\)-bridges.

For a two-terminal graph \(H\) we write \(H/st\) for the multigraph
obtained by identifying \(s\) and \(t\) into a single vertex \(z\) (an
\(st\) edge becomes a loop, which plays no role for connectivity and is
discarded). Thus
\(\mu_{H/st}(zw)=[sw\in E(H)]+[tw\in E(H)]\in\{0,1,2\}\) for
\(w\ne s,t\), and \(H/st\) is simple away from \(z\).

\subsection{\texorpdfstring{Balloon graphs and the candidate
\(G_{n,m}\)}{Balloon graphs and the candidate G\_\{n,m\}}}\label{balloon-graphs-and-the-candidate-g_nm}

\begin{definition}[Balloon graphs, after Romero 2025]\label{def:balloon}
For $n\ge 4$ and $n\le m\le\binom n2$ the balloon graph $B_{n,m}$ is defined by: if $m\ge\binom{n-1}2+2$, $B_{n,m}$ is $K_{n-1}$ together with one further vertex joined to $m-\binom{n-1}2$ vertices of the clique; $B_{4,4}$ is a triangle with a pendant vertex; otherwise $B_{n,m}$ is $B_{n-1,m-1}$ with a new pendant vertex attached to a vertex of minimum degree. The two-terminal balloon graph $G_{n,m}$ is $B_{n,m}$ with terminals at distance equal to the diameter of $B_{n,m}$.
\end{definition}

Romero (2025, Theorem 10) proves that \(G_{n,m}\) is unique up to
isomorphism and that the locally most split reliable graphs in
\(\mathcal T_{n,m}\) are exactly the two-terminal graphs with the same
split polynomial as \(G_{n,m}\) (``split-equivalent'' to it). We shall
only need the explicit form of \(G_{n,m}\) in the ranges that matter,
which is read off directly from the definition (with \(k=\binom n2-m\)):

\begin{itemize}
\tightlist
\item
  \(0\le k\le n-2\) (so \(m\ge\binom{n-1}2+1\)): \(G_{n,m}=K_n-K_{1,k}\)
  where the missing star is centred at \(s\) and \(t\) is one of its
  leaves (for \(k=0\) the terminals are any two vertices). Thus \(s\)
  has degree \(n-1-k\), every other vertex is adjacent to everything
  except possibly \(s\), and \(s,t\) are non-adjacent when \(k\ge 1\).
\item
  \(k=n-1\): \(s\) is a pendant vertex attached to \(u\), the other
  \(n-1\) vertices induce \(K_{n-1}\) minus one edge \(ux\), and
  \(t=x\).
\item
  \(k=n\): \(s\) is a pendant vertex attached to \(u\), the other
  \(n-1\) vertices induce \(K_{n-1}\) minus the two edges \(ux,uy\), and
  \(t=x\).
\item
  \(n+1\le k\le 2n-6\) (regime A1 below): \(s\) is pendant at \(u\); the
  other \(n-1\) vertices induce \(K_{n-1}-K_{1,r}\) with the missing
  star centred at \(u\), where \(r=k-n+2\in[3,n-4]\); and \(t\) is a
  leaf of that star (a non-neighbour of \(u\)).
\item
  \(2n-5\le k\le 3n-10\) (regime A2 below): \(s\) is the end of a
  pendant path \(s\!-\!w\!-\!u\) of length two; the other \(n-2\)
  vertices induce \(K_{n-2}-K_{1,r''}\) with the missing star centred at
  \(u\), where \(r''=k-2n+5\in[0,n-5]\); and \(t\) is a leaf of that
  star (any vertex of the clique other than \(u\) when \(r''=0\)).
\end{itemize}

These descriptions specify the representatives used below. The proof
uses equality of split polynomials, not uniqueness of all optimal
two-terminal graphs.

\subsection{\texorpdfstring{Behaviour near \(p=0\) and
\(p=1\)}{Behaviour near p=0 and p=1}}\label{behaviour-near-p0-and-p1}

The following elementary observation is Lemma 3 of Romero (2025),
attributed to Brown and McMullin.

\begin{lemma}\label{lem:ends}
Let $G,H\in\mathcal T_{n,m}$. If $N_i(G)=N_i(H)$ for all $i<i_0$ and $N_{i_0}(G)>N_{i_0}(H)$, then $\mathrm{SR}_G(p)>\mathrm{SR}_H(p)$ for all $p$ in some interval $(0,\delta)$. If $N_{m-j}(G)=N_{m-j}(H)$ for all $j<j_0$ and $N_{m-j_0}(G)>N_{m-j_0}(H)$, then $\mathrm{SR}_G(p)>\mathrm{SR}_H(p)$ for all $p$ in some interval $(1-\delta,1)$.
\end{lemma}

\begin{proof}
$\mathrm{SR}_G-\mathrm{SR}_H=\sum_i\big(N_i(G)-N_i(H)\big)p^i(1-p)^{m-i}$, and the first non-vanishing term dominates near the relevant endpoint.
\end{proof}

Consequently a UMSR graph must maximise \(N_{n-2}\) over
\(\mathcal T_{n,m}\) and must maximise the vector
\((N_{m-1},N_{m-2},\dots)\) in lexicographic order. Nonexistence is
proved by showing that these two requirements are incompatible.

\subsection{Split reliability as a difference of
reliabilities}\label{split-reliability-as-a-difference-of-reliabilities}

\begin{lemma}\label{lem:identity}
For every two-terminal graph $H$ and every $p$,
$$\mathrm{SR}_{s,t}(H;p)=\mathrm{Rel}(H/st;p)-\mathrm{Rel}(H;p),$$
where $\mathrm{Rel}(H/st;p)$ is the probability that the surviving edge set of $H$ is connected after $s$ and $t$ are identified. In particular $N_{n-2}(H)=\tau(H/st)$.
\end{lemma}

\begin{proof}
A spanning subgraph of $H$ is connected in $H/st$ if and only if it is connected in $H$, or it has exactly two components with $s$ and $t$ in different ones. The two events are disjoint, and the second is the split event. The last claim compares the coefficients of $p^{n-2}$ as $p\to0$: $H/st$ has $n-1$ vertices, so $\mathrm{Rel}(H/st;p)=\tau(H/st)p^{n-2}+O(p^{n-1})$, whereas $\mathrm{Rel}(H;p)=O(p^{n-1})$.
\end{proof}

The second statement is the classical fact that the number of spanning
\(2\)-forests of \(H\) separating \(s\) from \(t\) equals the number of
spanning trees of \(H\) with \(s\) and \(t\) identified. We shall use it
in both directions: in Section 3 to evaluate \(N_{n-2}\) of the
candidates and competitors through spanning-tree counts of dense graphs,
and in Section 4 to transfer the whole polynomial comparison to
all-terminal reliability.

\textbf{A three-vertex example and the loop convention.} Take a triangle
with terminals \(s,t\) and third vertex \(v\). A split state consists of
exactly one of \(sv,tv\): the edge \(st\) must be absent. Thus its split
reliability is \(2p(1-p)^2\). Identification leaves two parallel edges
from the merged vertex to \(v\), with reliability \(1-(1-p)^2=2p-p^2\);
the triangle itself has reliability \(3p^2-2p^3\). Subtracting gives
\(2p-4p^2+2p^3\), as required.

Discarding the identified loop does not change a probability, but it
does change edge-count coefficients. If \(X\) denotes the loopless
identification, with \(c_i=0\) outside its coefficient range, the exact
identity is \[N_i(H)=c_i(X)+\mathbf{1}_{st\in E(H)}c_{i-1}(X)-c_i(H).\]
The extra term counts states in which the discarded loop is selected. In
generating-polynomial form its contribution is multiplication by
\(1+x\). All coefficient comparisons below use a common total edge
count; silently subtracting coefficient vectors of different edge counts
is not valid.

\subsection{Spanning trees of nearly complete
graphs}\label{spanning-trees-of-nearly-complete-graphs}

\begin{lemma}\label{lem:trees}
Let $F$ be a graph on a subset of the vertices of $K_q$, with Laplacian $L_F$ (as a matrix indexed by the vertices of $F$). Then
$$\tau(K_q-F)=q^{\,q-2}\det\!\big(I-\tfrac1q L_F\big).$$
Consequently, for $a\ge 1$,
$$\tau(K_q-K_{1,a})=(q-a-1)(q-1)^{a-1}q^{\,q-a-2},\qquad
\tau\big(K_q-(K_{1,a}\cup K_2)\big)=(q-a-1)(q-1)^{a-1}(q-2)\,q^{\,q-a-3},$$
where in the second formula the star and the edge are vertex-disjoint.
\end{lemma}

\begin{proof}
$L(K_q-F)=qI-J-L_F$, and $J$ (the all-ones matrix) commutes with $L_F$ because $L_F$ has zero row sums. Hence the two matrices are simultaneously diagonalisable, the nonzero eigenvalues of $L(K_q-F)$ are $q-\mu$ for $\mu$ running over the Laplacian eigenvalues of $F$ (extended by zeros to all $q$ vertices) with one zero removed, and by the matrix-tree theorem $\tau(K_q-F)=\tfrac1q\prod(q-\mu)=\tfrac1{q^2}\det(qI-L_F)=q^{q-2}\det(I-\tfrac1qL_F)$. The Laplacian eigenvalues of $K_{1,a}$ are $a+1$, $1$ (with multiplicity $a-1$) and $0$; those of $K_2$ are $2$ and $0$; the displayed formulas follow.
\end{proof}

\subsection{A multigraph exchange
lemma}\label{a-multigraph-exchange-lemma}

For a loopless multigraph \(Y\) and distinct vertices \(x,y,z\) let
\(P_Y(xy\,|\,z)\) denote the probability (edges surviving independently
with probability \(p\)) that the surviving spanning subgraph has exactly
two connected components, one containing \(x\) and \(y\) and the other
containing \(z\); let \(P_Y(x\,|\,y)\) denote the probability of exactly
two components with \(x\) and \(y\) in different ones. For an edge
\(xy\) (possibly already present) write \(Y+xy\) for \(Y\) with one
additional \(xy\) edge. Adding the edge \(xy\) makes a surviving edge
set connected exactly when it was connected or had exactly two
components separating \(x\) from \(y\), so
\[\mathrm{Rel}(Y+xy;p)=\mathrm{Rel}(Y;p)+p\,P_Y(x\,|\,y).\qquad\text{(2.1)}\]

\begin{lemma}[Exchange lemma]\label{lem:exchange}
Let $Y$ be a loopless multigraph and $a,b,c$ distinct vertices. Suppose that
$$e_b(U)\le e_c(U)\qquad\text{for every }U\text{ with }a\in U\subseteq V(Y)\setminus\{b,c\}.$$
Then $P_Y(ab\,|\,c)\le P_Y(ac\,|\,b)$ for every $p\in[0,1]$; consequently $P_Y(a\,|\,c)\le P_Y(a\,|\,b)$ and $\mathrm{Rel}(Y+ab;p)\ge\mathrm{Rel}(Y+ac;p)$ for every $p$.
\end{lemma}

\begin{proof}
Let $E_0$ be the set of edges of $Y$ incident with neither $b$ nor $c$, and condition on the set $T$ of surviving edges of $E_0$. The components of $T$ partition $V'=V(Y)\setminus\{b,c\}$; call their vertex sets $U_1,\dots,U_r$ with $a\in U_1$. Given $T$, the events "$b$ has a surviving edge into $U_j$" (probability $\alpha_j=1-(1-p)^{e_b(U_j)}$), "$c$ has a surviving edge into $U_j$" (probability $\beta_j=1-(1-p)^{e_c(U_j)}$), for $j=1,\dots,r$, and "no $bc$ edge survives" (probability $(1-p)^{\mu(bc)}$) concern pairwise disjoint sets of edges and are therefore independent. The event $\{ab\,|\,c\}$ occurs exactly when no $bc$ edge survives, each $U_j$ is joined to exactly one of $b,c$, and $U_1$ is joined to $b$; hence
$$P(ab\,|\,c\mid T)=(1-p)^{\mu(bc)}\,\alpha_1(1-\beta_1)\prod_{j\ge2}\big(\alpha_j(1-\beta_j)+\beta_j(1-\alpha_j)\big),$$
and $P(ac\,|\,b\mid T)$ is the same expression with $\alpha_1(1-\beta_1)$ replaced by $\beta_1(1-\alpha_1)$. Since $\alpha_1(1-\beta_1)\le\beta_1(1-\alpha_1)$ is equivalent to $\alpha_1\le\beta_1$, that is to $e_b(U_1)\le e_c(U_1)$, which holds by hypothesis, the first claim follows by averaging over $T$. The second follows from $P_Y(a\,|\,c)=P_Y(ab\,|\,c)+P_Y(a\,|\,bc)$ and $P_Y(a\,|\,b)=P_Y(ac\,|\,b)+P_Y(a\,|\,bc)$, and the third from (2.1).
\end{proof}

This conditioning formulation belongs to the Kelmans/swing-surgery
tradition (Kelmans, 1981; Satyanarayana, Schoppmann and Suffel, 1992).
Brown, Colbourn and Devitt (1993) already extended reliability
transformations to multigraphs, as explicitly recorded by Csikvári
(2011, Section 2). We give a self-contained proof of the exact setwise
condition used here, without claiming priority for the multigraph
extension or a strict improvement over the 1993 hypotheses; the full
1993 text was not available for an exact comparison. The same
conditioning identifies the leading coefficient as \(p\to0\):

\begin{corollary}\label{cor:trees}
Under the hypothesis of Lemma \ref{lem:exchange},
$$\tau(Y+ab)-\tau(Y+ac)=\sum_{T}\big(e_c(U_1)-e_b(U_1)\big)\prod_{j\ge2}\big(e_b(U_j)+e_c(U_j)\big)\ \ge 0,$$
the sum over the spanning forests $T$ of $Y-\{b,c\}$, with $U_1,\dots,U_r$ the vertex sets of the trees of $T$ and $a\in U_1$. In particular, if $Y-\{b,c\}$ is connected and $\deg_Y(c)>\deg_Y(b)$ then $\tau(Y+ab)>\tau(Y+ac)$.
\end{corollary}

\begin{proof}
By (2.1), $\mathrm{Rel}(Y+ab)-\mathrm{Rel}(Y+ac)=p\big(P_Y(ac\,|\,b)-P_Y(ab\,|\,c)\big)$, and the difference of the two conditional probabilities in the proof of Lemma \ref{lem:exchange} is $(1-p)^{\mu(bc)}(\beta_1-\alpha_1)\prod_{j\ge2}(\cdots)$. As $p\to0$, a given edge set $T$ with $r$ components has probability $p^{|T|}(1+O(p))$, which is $p^{|V'|-r}(1+O(p))$ when $T$ is a forest and is of higher order otherwise; $\beta_1-\alpha_1=p\,(e_c(U_1)-e_b(U_1))+O(p^2)$; and each factor $\alpha_j(1-\beta_j)+\beta_j(1-\alpha_j)=p\,(e_b(U_j)+e_c(U_j))+O(p^2)$. Collecting terms, $\mathrm{Rel}(Y+ab)-\mathrm{Rel}(Y+ac)=p^{|V'|+1}\sum_T(\cdots)+O(p^{|V'|+2})$, and $|V'|+1=|V(Y)|-1$ is the exponent of $p$ carrying $\tau$ in $\mathrm{Rel}$. Each term of the sum is nonnegative by hypothesis; a spanning tree $T$ of $Y-\{b,c\}$ contributes $e_c(V')-e_b(V')=\deg_Y(c)-\deg_Y(b)$.
\end{proof}

\section{Nonexistence in the remaining intermediate-density
range}\label{nonexistence-in-the-remaining-intermediate-density-range}

We use Romero's (2025, Theorem 10) characterization of locally most
split reliable graphs: for the parameters considered here, every such
graph is split-equivalent to the canonical balloon candidate
\(G_{n,m}\). In particular, a UMSR graph \(H\) must satisfy
\[\mathrm{SR}_H=\mathrm{SR}_{G_{n,m}},\qquad N_{n-2}(H)=N_{n-2}(G_{n,m}).\]
Consequently it suffices to construct a connected simple two-terminal
graph with the same numbers of vertices and edges and a strictly larger
\(N_{n-2}\) coefficient. Such a graph beats \(G_{n,m}\) for all
sufficiently small positive \(p\). This argument concerns
split-equivalence, not uniqueness up to graph isomorphism.

\subsection{Regime A1: one pendant
terminal}\label{regime-a1-one-pendant-terminal}

\begin{proposition}\label{prop:A1comp}
Let $n\ge7$, $n+1\le k\le2n-6$, $m=\binom n2-k$, and $r=k-n+2$. No graph in $\mathcal T_{n,m}$ is UMSR.
\end{proposition}

\begin{proof}
Here $3\le r\le n-4$. The canonical candidate consists of a core $C=K_{n-1}-K_{1,r}$, with missing star centred at $u$, and a pendant terminal $s$ at $u$; the other terminal $x$ is a missing-star leaf. Identifying $s$ and $x$ restores $ux$, so
$$N_{n-2}(G_{n,m})=\tau(K_{n-1}-K_{1,r-1}).$$

Construct $H'$ with a pendant terminal $s$ at $u$, other terminal $x$, and core
$$H'-s=K_{n-1}-(K_{1,r-1}\cup K_2),$$
where the missing star is centred at $u$, contains $x$ as a leaf, and is vertex-disjoint from the additional missing edge $bc$. The construction uses $r+2\le n-2$ core vertices. The core induced by vertices other than $u$ is a clique minus $bc$ and is connected; $u$ has $n-1-r\ge3$ neighbours in it. Thus $H'$ is connected and simple. Its number of missing edges is $(n-2)+(r-1)+1=k$.

Identification restores $ux$. Applying the spanning-tree formulas of Lemma \ref{lem:trees},
$$\frac{N_{n-2}(H')}{N_{n-2}(G_{n,m})}
=\frac{\tau(K_{n-1}-(K_{1,r-2}\cup K_2))}{\tau(K_{n-1}-K_{1,r-1})}
=\frac{(n-r)(n-3)}{(n-1-r)(n-2)}>1,$$
since the numerator minus denominator of the last fraction is $r-2>0$. The candidate reduction above proves the assertion.
\end{proof}

\subsection{Regime A2: a pendant path of length
two}\label{regime-a2-a-pendant-path-of-length-two}

The preceding exchange lemma supplies a strict spanning-tree comparison
without requiring a closed formula for every missing-star size.

\begin{proposition}\label{prop:A2comp}
Let $n\ge7$, $2n-5\le k\le3n-10$, $m=\binom n2-k$, and $r=k-2n+5$. No graph in $\mathcal T_{n,m}$ is UMSR.
\end{proposition}

\begin{proof}
Put $q=n-2$, so $0\le r\le q-3$. Define a core $C$ on $q$ vertices as follows. For $r\ge1$, let $C=K_q-K_{1,r}$ with missing star centred at $u$ and choose a missing leaf $x$. For $r=0$, let $C=K_q$ and choose distinct vertices $u,x$. In particular, the zero-star case does not require a nonexistent missing leaf.

The canonical candidate is obtained by adjoining the path $s-w-u$ to $C$, with terminals $s,x$. Its identified graph $\Gamma=G_{n,m}/sx$ consists of $C$ and the vertex $w$ joined exactly to $u,x$.

Choose $\ell\in V(C)\setminus\{u,x\}$ adjacent to $u$ in $C$. There are $q-r-1\ge2$ choices when $r\ge1$, and $q-2\ge3$ when $r=0$. Choose any $a\in V(C)\setminus\{u,x,\ell\}$. The edge $a\ell$ belongs to $C$, whereas $aw$ is absent from $\Gamma$. Set
$$\Gamma'=\Gamma-a\ell+aw.$$
Form $H'$ by deleting $wx$ from $\Gamma'$ and adjoining a new pendant vertex $s$ at $w$, with terminals $s,x$. This is a simple graph with $n$ vertices and $\binom q2-r+2=m$ edges. It is connected: the subgraph of $C-a\ell$ induced by $V(C)\setminus\{u\}$ is a clique minus one edge on $q-1\ge4$ vertices, and $u$ retains its $q-1-r\ge2$ neighbours. Vertex $w$ is joined to $u$ and $a$, and $s$ to $w$. Identification of $s,x$ restores $wx$, giving $H'/sx=\Gamma'$.

To compare spanning-tree counts, put $Y=\Gamma-a\ell$. In $Y$, $w$ has precisely the neighbours $u,x$, while $\ell$ is adjacent to every vertex other than $w,a,\ell$. Therefore, for every $a\in U\subseteq V(Y)\setminus\{w,\ell\}$,
$$e_w(U)=|U\cap\{u,x\}|\le |U|-1=e_\ell(U).$$
Moreover, $Y-\{w,\ell\}=C-\ell$ is connected: its vertices other than $u$ form a clique, and deleting $\ell$ leaves at least $q-2-r\ge1$ neighbours of $u$. Finally,
$$\deg_Y(\ell)=q-2=n-4>2=\deg_Y(w).$$
The strict spanning-tree corollary \ref{cor:trees}, with $(b,c)=(w,\ell)$, now gives
$$N_{n-2}(H')=\tau(Y+aw)>\tau(Y+a\ell)=\tau(\Gamma)=N_{n-2}(G_{n,m}).$$
This covers $r=0$ as well as $r\ge1$, and the candidate reduction proves the proposition.
\end{proof}

\begin{theorem}\label{thm:A}
For $n\ge7$ and $\binom{n-3}2+4\le m\le\binom n2-n-1$, no graph in $\mathcal T_{n,m}$ is UMSR. Together with Romero (2025, Theorem 12), this gives nonexistence for every $n\le m\le\binom n2-n-1$.
\end{theorem}

\begin{proof}
The two propositions cover all missing-edge counts $n+1\le k\le3n-10$, precisely the displayed range of $m$. Romero's Theorem 12 covers $n\le m\le\binom{n-3}2+3$.
\end{proof}

\section{\texorpdfstring{Existence for \(m\ge\binom n2-n+2\), and the
two boundary
values}{Existence for m\textbackslash ge\textbackslash binom n2-n+2, and the two boundary values}}\label{existence-for-mgebinom-n2-n2-and-the-two-boundary-values}

\subsection{The identified graph of a two-terminal
graph}\label{the-identified-graph-of-a-two-terminal-graph}

Fix \(n\ge 4\), \(1\le k\le n\) and \(m=\binom n2-k\). Let
\(d=n-1-k\ge-1\). Define the loopless multigraph \(X^{(k)}\) on \(n-1\)
vertices as follows: for \(k\le n-1\), \(X^{(k)}\) is \(K_{n-1}\) with
\(d\) additional edges parallel to \(d\) distinct edges at one fixed
vertex \(z\) (so \(X^{(n-1)}=K_{n-1}\)); and \(X^{(n)}=K_{n-1}-e\). From
the explicit description of \(G_{n,m}\) in the balloon-graph definition
one checks directly that \(G_{n,m}/st\cong X^{(k)}\): for \(k\le n-2\)
the merged vertex is joined twice to each of the \(n-1-k\) neighbours of
\(s\) (all of which are neighbours of \(t\)) and once to every other
vertex; for \(k=n-1\) the neighbourhoods of \(s\) and \(t\) are disjoint
and together cover the other vertices; for \(k=n\) the same holds except
that the edge \(uy\) is missing from the clique.

\begin{lemma}\label{lem:identified}
For every $H\in\mathcal T_{n,m}$ and every $p\in[0,1]$, $\ \mathrm{Rel}(H/st;p)\le\mathrm{Rel}(X^{(k)};p)$.
\end{lemma}

\begin{proof}
Let $W=V(H)\setminus\{s,t\}$ and let $\mathfrak X$ be the family of loopless multigraphs $X$ on $W\cup\{z\}$ with exactly $m$ edges such that $X[W]$ is simple and $\mu(zw)\le2$ for all $w\in W$. If $st\notin E(H)$ then $H/st\in\mathfrak X$. If $st\in E(H)$, then $H/st$ (with its loop discarded) has $m-1$ edges, and since $k\ge1$ some edge of $K_n$ other than $st$ is missing from $H$; adding the corresponding edge (a $W$-edge, or one more copy of a $z$-edge) produces a member $X$ of $\mathfrak X$ with $\mathrm{Rel}(H/st)\le\mathrm{Rel}(X)$ by monotonicity. So it suffices to prove $\mathrm{Rel}(X)\le\mathrm{Rel}(X^{(k)})$ for $X\in\mathfrak X$.

Describe $X\in\mathfrak X$ by the set $F_W$ of pairs of $W$ not joined in $X$ ($f_W=|F_W|$), the set $Z_0=\{w:\mu(zw)=0\}$ and the set $D_2=\{w:\mu(zw)=2\}$. Counting edges,
$$|D_2|=d+f_W+|Z_0|.\qquad\text{(4.2)}$$
Call $X$ *terminal* if $f_W=0$ and $Z_0=\emptyset$ when $d\ge0$, and if $f_W+|Z_0|=1$ and $D_2=\emptyset$ when $d=-1$. A terminal $X$ is isomorphic to $X^{(k)}$: for $d\ge0$ it is $K_{n-1}$ with exactly $d$ doubled edges at $z$ by (4.2); for $d=-1$ it is $K_{n-1}$ minus one edge. We show that every non-terminal $X\in\mathfrak X$ admits a move to some $X'\in\mathfrak X$ with $\mathrm{Rel}(X')\ge\mathrm{Rel}(X)$ that strictly decreases $\Phi(X)=2f_W+|Z_0|$; since $\Phi\ge0$ the process ends at a terminal graph, proving the lemma.

*Move M1 (fill a missing $W$-edge).* Let $xy\in F_W$ with $\mu(zx)\ge1$ and $|Z_0\cap N_W(y)|\le\mu(zx)-1$, where $N_W(y)$ is the set of $W$-neighbours of $y$ in $X$. Let $Y$ be $X$ with one copy of $zx$ removed and $X'=Y+xy$. For $U$ with $x\in U\subseteq W\setminus\{y\}$,
$$e_z(U)-e_y(U)=\big(\mu_X(zx)-1\big)+\sum_{u\in U\setminus\{x\}}\big(\mu_X(zu)-[uy\in E(X)]\big)\ \ge\ \mu_X(zx)-1-|Z_0\cap N_W(y)|\ \ge 0,$$
so Lemma \ref{lem:exchange} with $(a,b,c)=(x,y,z)$ gives $\mathrm{Rel}(X')=\mathrm{Rel}(Y+xy)\ge\mathrm{Rel}(Y+xz)=\mathrm{Rel}(X)$. The move removes $xy$ from $F_W$ and adds $x$ to $Z_0$ only if $\mu_X(zx)=1$, so $\Phi$ decreases by at least $1$.

*Move M3 (fill a missing $z$-edge).* Let $w\in Z_0$ and $v\in D_2$ with at most one $F_W$-neighbour of $v$ outside $N_{F_W}(w)\cup\{w\}$, where $N_{F_W}$ denotes neighbourhoods in the graph $F_W$. Let $Y$ be $X$ with one copy of $zv$ removed and $X'=Y+zw$. For $U$ with $z\in U\subseteq V(X)\setminus\{v,w\}$, the contribution of $z$ to $e_v(U)-e_w(U)$ is $\mu_Y(vz)-\mu_Y(wz)=1$, and each $u\in U\setminus\{z\}$ contributes $[uv\in E]-[uw\in E]=[uw\in F_W]-[uv\in F_W]\ge-1$, with $-1$ exactly when $uv\in F_W$ and $uw\notin F_W$. Hence $e_v(U)-e_w(U)\ge1-1=0$, and Lemma \ref{lem:exchange} with $(a,b,c)=(z,w,v)$ gives $\mathrm{Rel}(X')\ge\mathrm{Rel}(X)$. The move removes $w$ from $Z_0$ and $v$ from $D_2$, so $\Phi$ decreases by $1$.

*A move always exists.* Suppose $X$ is non-terminal. If $Z_0=\emptyset$ then $f_W\ge1$ (for $d\ge0$) or $f_W\ge2$ (for $d=-1$), and M1 applies to any $xy\in F_W$ because its hypothesis reduces to $\mu(zx)\ge1$. If $Z_0\ne\emptyset$, then $D_2\ne\emptyset$ by (4.2) (for $d=-1$ non-terminality gives $f_W+|Z_0|\ge2$). If some $v\in D_2$ has $F_W$-degree at most $1$, M3 applies with any $w\in Z_0$. Otherwise every $v\in D_2$ has $F_W$-degree at least $2$, so $2f_W\ge\sum_{v\in D_2}\deg_{F_W}(v)\ge2|D_2|=2(d+f_W+|Z_0|)$, forcing $d=-1$, $|Z_0|=1$, and equality throughout: every edge of $F_W$ joins two vertices of $D_2$ of $F_W$-degree exactly $2$. Then $f_W\ge1$ by non-terminality, and for any $xy\in F_W$ we have $\mu(zx)=2$ and $|Z_0\cap N_W(y)|\le1$, so M1 applies.
\end{proof}

\subsection{Proof of parts (iii) and
(iv)}\label{proof-of-parts-iii-and-iv}

\begin{theorem}\label{thm:B}
Let $n\ge 4$, $0\le k\le n-2$ and $m=\binom n2-k$. Then $G_{n,m}$ is uniformly most split reliable in $\mathcal T_{n,m}$.
\end{theorem}

\begin{proof}
For $k=0$ there is nothing to prove. For $k\ge1$, by Lemma \ref{lem:identity}, Lemma \ref{lem:identified}, and the isomorphism $G_{n,m}/st\cong X^{(k)}$,
$$\mathrm{SR}_H(p)=\mathrm{Rel}(H/st;p)-\mathrm{Rel}(H;p)\le\mathrm{Rel}(G_{n,m}/st;p)-\mathrm{Rel}(H;p).$$
As an unlabelled graph $G_{n,m}$ is the balloon graph $K_{n-1}$ plus a vertex of degree $n-1-k\ge1$, and $m\ge\binom{n-1}2+1$; by the theorem of Petingi, Saccoman and Schoppmann (1996) this graph is uniformly least reliable among connected graphs with $n$ vertices and $m$ edges, so $\mathrm{Rel}(H;p)\ge\mathrm{Rel}(G_{n,m};p)$ for all $p$. Hence $\mathrm{SR}_H(p)\le\mathrm{SR}_{G_{n,m}}(p)$ for all $p$, which is what the definition requires.
\end{proof}

\begin{theorem}\label{thm:Bprime}
Let $n\ge 7$ and $k\in\{n-1,n\}$, $m=\binom n2-k$. If $B_{n,m}$ is uniformly least reliable among connected graphs with $n$ vertices and $m$ edges, then $G_{n,m}$ is uniformly most split reliable in $\mathcal T_{n,m}$. The hypothesis holds for $7\le n\le 51$ (Corollary \ref{cor:n41}).
\end{theorem}

\begin{proof}
The displayed inequality in the proof of Theorem \ref{thm:B} holds verbatim with Lemma \ref{lem:identified} for $k\in\{n-1,n\}$, and $G_{n,m}$ is $B_{n,m}$ as an unlabelled graph; the hypothesis supplies $\mathrm{Rel}(H)\ge\mathrm{Rel}(G_{n,m})$.
\end{proof}

The hypothesis of Theorem \ref{thm:Bprime} is only needed for graphs of
minimum degree at least \(2\): if \(H\) has a pendant vertex then
\(\mathrm{Rel}(H)=p\,\mathrm{Rel}(H')\) where \(H'\) has \(n-1\)
vertices and \(m-1\ge\binom{n-2}2+1\) edges, and the
Petingi--Saccoman--Schoppmann theorem applied to \(H'\) gives
\(\mathrm{Rel}(H)\ge\mathrm{Rel}(B_{n,m})\). The rest of this section
reduces the hypothesis, for all \(n\), to an explicit family of
inequalities between reliabilities of cones over a complete graph, and
verifies that inequality exactly for \(n\le 51\).

\subsection{Reduction of the hypothesis to an inequality between two
cones}\label{reduction-of-the-hypothesis-to-an-inequality-between-two-cones}

For \(N\ge3\) and \(0\le s\le N-1\) let \(C_N(s)=K_N-K_{1,N-1-s}\)
denote the \emph{cone graph}: a vertex \(u\) joined to \(s\) vertices of
a clique \(K_{N-1}\). Thus
\(K_{N}-K_{1,1+\varepsilon}=C_N(N-2-\varepsilon)\), and \(B_{n,m}\) for
\(m=\binom{n-1}2-\varepsilon\) is \(C_{n-1}(n-3-\varepsilon)\) with a
pendant vertex attached to a clique vertex. Write
\(f_N(s;p)=\mathrm{Rel}(C_N(s);p)\) and \(q=1-p\).

\begin{definition}[Inequality $(\star)_{N,\varepsilon}$]\label{def:star}
Let $N\ge 7$ and $\varepsilon\in\{0,1\}$, and put $h=N-2-\varepsilon$. We say that $(\star)_{N,\varepsilon}$ holds if for every integer $d$ with $1\le d\le N-5-\varepsilon$ and every $p\in[0,1]$,
$$(1-q^{d+1})\,f_N(h-d;p)\;\ge\;p\,f_N(h;p).\qquad\text{(4.3)}$$
\end{definition}

Both sides of (4.3) are reliabilities of multigraphs on \(N+1\) vertices
with \(\binom{N-1}2+h+1\) edges: the right side is the balloon
(\(C_N(h)\) with a pendant vertex \(w\) attached to a clique vertex
\(x\)), the left side is \(C_N(h-d)\) with \(w\) attached to \(x\) by
\(d+1\) parallel edges. So \((\star)\) says that taking \(d\) of the
cone edges away from \(u\) and giving them to the pendant vertex, as
parallel copies of its one edge, does not decrease reliability.
Expanding \(1-q^{d+1}=\sum_{i\ge1}\binom{d+1}i p^iq^{d+1-i}\), (4.3) is
implied by the coefficient-wise statement
\[\sum_{i=1}^{d+1}\binom{d+1}{i}\,c_{j-i}\bigl(C_N(h-d)\bigr)\;\ge\;c_{j-1}\bigl(C_N(h)\bigr)\qquad\text{for all }j,\qquad\text{(4.4)}\]
because the basis \(p^jq^{M-j}\) is nonnegative on \([0,1]\). Here
\(c_j(G)\) is the number of connected spanning subgraphs of \(G\) with
\(j\) edges.

The reduction uses Kelmans's transformation. For two vertices \(u,v\) of
a graph \(G\), the \emph{Kelmans transformation} with beneficiary \(u\)
and co-beneficiary \(v\) deletes all edges \(vw\) with
\(w\in N(v)\setminus(N(u)\cup\{u\})\) and adds the edges \(uw\) for the
same \(w\); it preserves the number of edges. Kelmans (1981) proved that
it never increases all-terminal reliability (see Csikvári, 2011, Theorem
2.2, with \(k=1\) there), and Csikvári (2011, Theorem 2.6) proved that
every graph can be carried by a sequence of Kelmans transformations to a
threshold graph with the same number of components. Recall that a graph
is \emph{threshold} if any two vertices \(u,v\) have comparable sets
\(N(u)\setminus\{v\}\) and \(N(v)\setminus\{u\}\); equivalently it is
obtained from a single vertex by repeatedly adding an isolated or a
dominating vertex. The complement of a threshold graph is threshold, and
a threshold graph without isolated vertices has a vertex adjacent to all
others.

\begin{theorem}\label{thm:reduction}
Let $\varepsilon\in\{0,1\}$, $n\ge 8+3\varepsilon$, $m=\binom{n-1}2-\varepsilon$ and $N=n-1$. If $(\star)_{N,\varepsilon}$ holds, then every connected graph $H$ with $n$ vertices and $m$ edges satisfies
$$\mathrm{Rel}(H;p)\;\ge\;\mathrm{Rel}(B_{n,m};p)=p\,\mathrm{Rel}(K_{n-1}-K_{1,1+\varepsilon};p)\qquad(0\le p\le1),$$
i.e. $B_{n,m}$ is uniformly least reliable, and consequently $G_{n,m}$ is uniformly most split reliable in $\mathcal T_{n,m}$.
\end{theorem}

\begin{proof}
By the results of Kelmans and Csikvári just quoted there is a connected threshold graph $H_{\mathrm{tr}}$ with $n$ vertices and $m$ edges and $\mathrm{Rel}(H_{\mathrm{tr}})\le\mathrm{Rel}(H)$ on $[0,1]$, so we may assume that $H$ itself is threshold. Let $F$ be the complement of $H$; it is a threshold graph with $n-1+\varepsilon$ edges. Since $H$ is connected and threshold, its last vertex in a threshold construction sequence is dominating, so $F$ has an isolated vertex. Let $t$ be the number of non-isolated vertices of $F$ and $F_0$ the subgraph they induce; $F_0$ is threshold without isolated vertices, so it has a vertex $w$ adjacent in $F$ to the other $t-1$ vertices of $F_0$. In $H$ the neighbours of $w$ are exactly the $n-t$ vertices outside $F_0$, and each of them is isolated in $F$, hence dominating in $H$. Because $\binom t2\ge e(F_0)=n-1+\varepsilon$ and $n\ge8+3\varepsilon$, we have $t\ge5+\varepsilon$ (for $\varepsilon=0$, $t\le4$ would give $n-1\le6$; for $\varepsilon=1$, $t\le5$ would give $n\le10$). Also $t\le n-1$.

If $t=n-1$ then $w$ is pendant, $\mathrm{Rel}(H)=p\,\mathrm{Rel}(H-w)$, and $H-w$ is a connected graph with $n-1$ vertices and $\binom{n-1}2-1-\varepsilon\ge\binom{n-2}2+1$ edges, so by the Petingi–Saccoman–Schoppmann theorem $\mathrm{Rel}(H-w)\ge\mathrm{Rel}(K_{n-1}-K_{1,1+\varepsilon})$ and we are done.

Otherwise put $d=n-t-1$, so that $w$ has degree $d+1\ge2$ and $1\le d\le n-6-\varepsilon=N-5-\varepsilon$. The graph $H$ is connected whenever $H-w$ is connected and at least one of the $d+1$ edges at $w$ is present, and these two events concern disjoint edge sets; hence
$$\mathrm{Rel}(H)\ \ge\ (1-q^{d+1})\,\mathrm{Rel}(H-w).$$
The graph $H-w$ is connected (any neighbour of $w$ dominates it), has $N$ vertices, and has $(n-1+\varepsilon)-(t-1)=d+1+\varepsilon$ missing edges, i.e. $\binom N2-(d+1+\varepsilon)\ge\binom{N-1}2+1$ edges because $d\le N-3-\varepsilon$. By the Petingi–Saccoman–Schoppmann theorem,
$$\mathrm{Rel}(H-w)\ \ge\ \mathrm{Rel}\bigl(K_N-K_{1,d+1+\varepsilon}\bigr)=f_N(N-2-\varepsilon-d)=f_N(h-d).$$
Combining the two displays with $(\star)_{N,\varepsilon}$ gives $\mathrm{Rel}(H)\ge(1-q^{d+1})f_N(h-d)\ge p\,f_N(h)=p\,\mathrm{Rel}(K_{n-1}-K_{1,1+\varepsilon})$. The last assertion is Theorem \ref{thm:Bprime}.
\end{proof}

The next proposition separates finite certification from an analytic
endpoint comparison. The first excluded value has five small-order
exceptions; it does not fail uniformly over all orders.

\begin{proposition}\label{prop:star}
Let $\varepsilon\in\{0,1\}$, $N\ge7$ and $h=N-2-\varepsilon$.
\begin{itemize}
\item[(a)] For $7\le N\le 50$, the coefficient-wise inequality (4.4) holds for every $1\le d\le N-5-\varepsilon$ and every $j$; in particular $(\star)_{N,\varepsilon}$ holds for $7\le N\le50$.
\item[(b)] For every $N\ge7$ and $1\le d\le N-5-\varepsilon$, (4.3) holds to leading order as $p\to0$: $(d+1)\,\tau(C_N(h-d))\ge\tau(C_N(h))$.
\item[(c)] At the first excluded value $d=N-4-\varepsilon=h-2$, (4.4) holds in the five cases $(N,\varepsilon)=(7,0),(7,1),(8,1),(9,1),(10,1)$. It fails in the other cases with $7\le N\le50$. For every $N\ge11$, the analytic obstruction is already in the leading coefficient: $(d+1)\tau(C_N(h-d))<\tau(C_N(h))$. Thus the proposed range cannot be enlarged uniformly in order.
\end{itemize}
\end{proposition}

\begin{proof}
(a) is a finite computation in exact integer arithmetic, described in Section 5.

(b) By Lemma \ref{lem:trees}, $\tau(C_N(s))=\tau(K_N-K_{1,N-1-s})=s\,(N-1)^{N-2-s}N^{s-1}$, so
$$\frac{(d+1)\,\tau(C_N(h-d))}{\tau(C_N(h))}=\frac{(d+1)(h-d)}{h}\Bigl(1-\frac1N\Bigr)^{d}.$$
Put $s=h-d\ge3$. The first factor $(d+1)s/(s+d)$ is increasing in $s$, so it is at least $3(d+1)/(d+3)$, and $(1-1/N)^d\ge e^{-d/(N-1)}\ge e^{-d/(d+4)}$ because $N-1\ge h+1\ge d+4$. The function $\varphi(d)=\log\frac{3(d+1)}{d+3}-\frac{d}{d+4}$ satisfies $\varphi(1)=\log1.5-0.2>0$, $\varphi(d)\to\log3-1>0$ as $d\to\infty$, and $\varphi'(d)=\frac{2}{(d+1)(d+3)}-\frac{4}{(d+4)^2}$ changes sign only once (from positive to negative), so $\varphi>0$ for all $d\ge1$; hence the ratio is at least $1$.

(c) For $d=h-2$ the ratio in (b) equals $r(N,\varepsilon)=\frac{2(h-1)}{h}(1-\frac1N)^{h-2}$. Since $h\ge N-3$, $r<2(1-\frac1N)^{N-5}\le 2e^{-(N-5)/N}$, which is below $1$ for $N\ge17$; for $11\le N\le16$ direct evaluation gives $r\le0.988$ (the values at $N=11$ are $0.912$ for $\varepsilon=0$ and $0.988$ for $\varepsilon=1$). The excluded value is tested separately from (a). Exact coefficient checks give the five stated exceptions and failures in the remaining cases; these checks are retained separately in the evidence package. A negative leading coefficient proves failure near $p=0$, whereas an arbitrary negative coefficient alone would not suffice to refute the probability inequality.
\end{proof}

\begin{corollary}\label{cor:n41}
For $7\le n\le 51$ and $m\in\{\binom{n-1}2,\binom{n-1}2-1\}$ the balloon graph $B_{n,m}$ is uniformly least reliable among connected graphs with $n$ vertices and $m$ edges, and $G_{n,m}$ is uniformly most split reliable in $\mathcal T_{n,m}$. For every $n\ge 11$ the same conclusions follow from $(\star)_{n-1,\varepsilon}$ alone.
\end{corollary}

\begin{proof}
The required base cases $7\le n\le10$ are covered by exact all-terminal coefficient comparisons in Section 5. For $8\le n\le51$ with $\varepsilon=0$ and $11\le n\le51$ with $\varepsilon=1$ it is Theorem \ref{thm:reduction} together with Proposition \ref{prop:star}(a), applied with $N=n-1\le50$.
\end{proof}

Theorem \ref{thm:reduction} replaces a statement about all connected
graphs with \(n\) vertices and \(m\) edges by a statement about one
explicit pair of graphs. The two sides of (4.3) have equal endpoint
values, not equal first-order expansions. As \(q\to0\), their difference
is \(q-q^{d+1}-q^{h-d}+q^h+O(q^{N-2})>0\), and as \(p\to0\) it is
governed by the spanning-tree ratio of Proposition \ref{prop:star}(b),
which exceeds \(1\) by a bounded factor. We have not found a proof of
(4.3) for all \(N\); an equivalent probabilistic form is that the
probability that the \(d\) extra cone edges are needed to connect
\(C_N(h)\) is at most \((q+q^2+\dots+q^d)\,f_N(h-d)\).

\emph{Proof of Theorem \ref{thm:main}.} Part (i) is Proposition 2.1 of
Brown and McMullin (2023). Part (ii) is Theorem \ref{thm:A}. Part (iii)
is Theorem \ref{thm:B} and part (iv) is Theorem \ref{thm:Bprime}; its
proviso holds for \(7\le n\le 51\) by Corollary \ref{cor:n41}, which
gives the final sentence for these \(n\). \(\square\)

\subsection{Analytic intervals valid at every
order}\label{analytic-intervals-valid-at-every-order}

The finite coefficient certificates and the following analytic
comparison provide different kinds of coverage. The certificates cover
the whole probability interval at bounded orders. The analytic
comparison applies at every order, but leaves an intermediate
probability interval unresolved.

\begin{theorem}[Cone comparison on explicit probability intervals]\label{thm:coneintervals}
Let $N\ge7$, $3\le s<h\le N-1$, and $d=h-s$. Then
$$
(1-q^{d+1})f_N(s;p)\ge p f_N(h;p),\qquad q=1-p,
$$
for every $p\in[1/2,1]$, with strict inequality for $1/2\le p<1$.

In the particular range $h=N-2-\varepsilon$, $\varepsilon\in\{0,1\}$ and $1\le d\le N-5-\varepsilon$, put
$$
R=\frac{(d+1)(h-d)}{h}\left(\frac{N-1}{N}\right)^d,
\qquad L=\binom{N-1}{2}-1-\varepsilon.
$$
Then $R>1$, and the same comparison holds on
$$0\le p\le1-R^{-1/L}.$$
\end{theorem}

\begin{proof}
Write $M=N-1$, $r_M(p)=\operatorname{Rel}(K_M;p)$ and $A=1+q+\cdots+q^d$. A sufficient event for the smaller cone to be connected is that the surviving clique is connected and at least one of its $s$ cone edges survives. These events use disjoint edge sets. Conversely, connectivity of the larger cone requires at least one of its $h$ cone edges to survive. Hence
$$f_N(s;p)\ge(1-q^s)r_M(p),\qquad f_N(h;p)\le1-q^h.$$
For $0<q<1$ it therefore suffices to prove
$$
1-r_M(p)\le D:=1-\frac{1-q^{s+d}}{A(1-q^s)}
=q\frac{1-q^d}{1-q^{d+1}}\frac{1-q^{s-1}}{1-q^s}.
$$
For fixed $q$, the sequence $(1-q^a)/(1-q^{a+1})$ increases with the positive integer $a$: the difference of adjacent cross products is $q^a(1-q)^2>0$. Since $d\ge1$ and $s-1\ge2$,
$$D\ge q\frac{1-q}{1-q^2}\frac{1-q^2}{1-q^3}
=\frac{q}{1+q+q^2}.$$

Every disconnected clique state has a nontrivial cut all of whose edges have failed. Counting each unordered cut once, for $M=6$ the union bound gives
$$1-r_6(p)\le6q^5+15q^8+10q^9.$$
If $0<q\le1/2$, division by $q$ bounds this by
$$\frac6{16}+\frac{15}{128}+\frac{10}{256}
=\frac{17}{32}<\frac47\le\frac1{1+q+q^2}.$$
For $M\ge7$, singleton cuts contribute $Mq^{M-1}$; each remaining cut has at least $2(M-2)$ crossing edges, and there are fewer than $2^{M-1}$ such cuts. Thus
$$
\frac{1-r_M(p)}q
\le Mq^{M-2}+2^{M-1}q^{2M-5}
\le\frac{M}{2^{M-2}}+2^{4-M}
\le\frac7{32}+\frac18=\frac{11}{32}<\frac47.
$$
Both terms in the penultimate expression decrease with $M\ge7$. This proves the strict sufficient inequality for $0<q\le1/2$. At $p=1$ both sides of the cone comparison equal $1$.

For the lower interval, the smaller cone has $E_s=\binom M2+s$ edges. States consisting of exactly a spanning tree give the lower bound
$$f_N(s;p)\ge\tau(C_N(s))p^{N-1}q^{E_s-N+1}.$$
The union bound over spanning trees gives
$$f_N(h;p)\le\tau(C_N(h))p^{N-1}.$$
Together with $1-q^{d+1}\ge(d+1)pq^d$, these imply the desired comparison whenever $Rq^L\ge1$, since
$$
R=\frac{(d+1)\tau(C_N(s))}{\tau(C_N(h))},
\qquad d+E_s-N+1=\binom M2-1-\varepsilon=L.
$$
The cone tree formula yields the stated expression for $R$.

To verify that the lower interval has positive length, use $s\ge3$ and $N-1\ge d+4$ to obtain
$$
\log R\ge\log\frac{3(d+1)}{d+3}-\frac d{d+4}=:\varphi(d),
$$
where $\log(1-1/N)\ge-1/(N-1)$ was used. The derivative
$$\varphi'(d)=\frac2{(d+1)(d+3)}-\frac4{(d+4)^2}$$
has the sign of $10-d^2$, so $\varphi$ has only one maximum on $[1,\infty)$. Its value at $1$ is $\log(3/2)-1/5>0$ and its limit at infinity is $\log3-1>0$. Therefore $R>1$. The condition $Rq^L\ge1$ is exactly $p\le1-R^{-1/L}$; at $p=0$ the comparison holds directly.
\end{proof}

In particular, the threshold-graph reduction applies pointwise: for
\(\varepsilon\in\{0,1\}\), \(n\ge8+3\varepsilon\) and
\(m=\binom{n-1}{2}-\varepsilon\), every connected graph \(H\) with \(n\)
vertices and \(m\) edges satisfies
\[\operatorname{Rel}(H;p)\ge\operatorname{Rel}(B_{n,m};p)\qquad(1/2\le p\le1).\]
The same implication holds on the common lower interval obtained by
taking the minimum of \(1-R^{-1/L}\) over the finitely many required
values of \(d\) for \(N=n-1\). These are pointwise interval conclusions,
not uniform reliability on \([0,1]\). The intermediate interval remains
unresolved at unrestricted orders; the full-interval finite theorem
through \(n=51\) remains a separate result.

\section{Exact computations and their
scope}\label{exact-computations-and-their-scope}

The universal nonexistence and dense-existence arguments above are
written proofs. The finite boundary conclusion additionally depends on
exact cone coefficients and a small number of all-terminal base cases.
The other enumerations are corroboration, not substitutes for those
arguments. None of these programs is a proof-assistant formalisation of
the mathematics.

\subsection{Certificate semantics}\label{certificate-semantics}

For each graph, \(c_i\) counts connected spanning edge sets of size
\(i\), and \(N_i\) counts the corresponding split states. The
connected-subgraph generating polynomial on an induced vertex set \(S\)
is obtained by subtracting states whose component containing the least
vertex is a proper subset \(T\): \[c_S(x)=(1+x)^{e(S)}-
\sum_{\substack{T\subsetneq S\\\min S\in T}}
c_T(x)(1+x)^{e(S\setminus T)}.\] Edges between \(T\) and its complement
are absent. Split coefficients then come from
\[\sum_{\substack{A\subseteq V\\s\in A,\ t\notin A}}
c_A(x)c_{V\setminus A}(x).\] Counting every terminal pair does not
require treating isomorphic copies as different scientific objects. The
graph6 corpora contain unlabelled graph representatives; all their
unordered terminal pairs are checked when a full split-vector run is
specified.

A coefficientwise inequality at a common total edge count implies a
probability inequality, because \(p^i(1-p)^{m-i}\ge0\). It is
sufficient, not necessary. Nonexistence is certified differently: the
graphs maximising the lowest split coefficient and those maximising the
reversed coefficient vector lexicographically have disjoint sets of
coefficient vectors. The endpoint lemma then rules out a uniform
optimiser. Failure of a coefficientwise comparison alone is not such a
nonexistence certificate.

\subsection{Cone coefficients}\label{cone-coefficients}

Write \(M=N-1\). If \(a\) of the \(s\) cone edges survive, the surviving
clique must have every component meeting their \(a\) endpoints. Let
\(\gamma_M(a;x)\) count those clique edge sets. Removing the union of
components avoiding that fixed root set gives
\[\gamma_M(a;x)=(1+x)^{\binom M2}
-\sum_{u=1}^{M-a}\binom{M-a}{u}(1+x)^{\binom u2}\gamma_{M-u}(a;x),\] and
hence \[\sum_j c_j(C_N(s))x^j=\sum_{a=1}^s\binom sa x^a\gamma_M(a;x).\]
The fast implementation uses packed integer polynomials with an explicit
digit bound; an alternative implementation groups the generic subset
recursion by symmetry orbits. Ordinary-list arithmetic and direct
edge-subset checks provide smaller, transparent comparisons. Packing
arithmetic is an implementation choice, not a strengthening of the
claim.

\subsection{Recorded coverage}\label{recorded-coverage}

The incoming reviewer evidence records the following executed checks.
These are supplied review provenance, not executions newly attributed to
this revision. The release's replay receipt separately identifies
executions performed on the revised package and their exact input
hashes.

{\def\LTcaptype{none} % do not increment counter
\begin{longtable}[]{@{}
  >{\raggedright\arraybackslash}p{(\linewidth - 4\tabcolsep) * \real{0.3333}}
  >{\raggedright\arraybackslash}p{(\linewidth - 4\tabcolsep) * \real{0.3333}}
  >{\raggedright\arraybackslash}p{(\linewidth - 4\tabcolsep) * \real{0.3333}}@{}}
\toprule\noalign{}
\begin{minipage}[b]{\linewidth}\raggedright
Check
\end{minipage} & \begin{minipage}[b]{\linewidth}\raggedright
Exact scope
\end{minipage} & \begin{minipage}[b]{\linewidth}\raggedright
Role
\end{minipage} \\
\midrule\noalign{}
\endhead
\bottomrule\noalign{}
\endlastfoot
Required cone inequalities & \(7\le N\le50\), both \(\varepsilon\),
every required \(d\): 2,024 triples and 1,354,694 coefficient
comparisons & Finite boundary theorem \\
Boundary all-terminal minima & Both boundary edge counts at
\(n=7,8,9,10\); 5,422 graphs & Base cases for the reduction \\
Rebuilt C split vectors & All connected graphs at \(n=5,6,7,8\), and the
supplied dense \(n=10\) corpus; 16,678 graphs, 536,954 terminal pairs &
Implementation replay \\
Direct edge-subset comparison & All 133 connected graphs at \(n=5,6\);
1,890 terminal pairs & Independent algorithm check within supplied
review \\
Root-set and generic cone checks & 21 root-set checks through clique
order 6; 39 cone checks at orders 5--10 & Formula and implementation
check \\
Additional dense \(n=9\) replay & 1,370 graphs, 49,320 terminal pairs &
Supplementary corroboration \\
Corpus regeneration & Nine supplied graph6 corpora & Input enumeration
check \\
First excluded cone value & 88 cases: 83 leading-coefficient failures
and five full coefficientwise passes & Correction of false sharpness
assertion \\
\end{longtable}
}

The original identity/identified-graph logs cover \textbf{3,640 terminal
pairs from 140 graph representatives}: 40 dense graphs of order 7 and
100 of order 8. They do not cover every two-terminal graph through order
8. For adjacent terminals the loop-padded coefficient identity from the
preliminaries is required.

The original fast cone logs reach \(N=50\). Complete two-\(\varepsilon\)
coverage of the slower implementation ends at \(N=30\), with only
\(\varepsilon=0\) recorded at \(N=31\). Its generic validation covers
\(N=5,6,7\); the faster implementation's validation covers
\(N=5,\ldots,10\). We do not describe these as two complete independent
runs through \(N=50\).

Original summaries also report full split enumeration at \(n=9\) and
dense \(n=11,12\). Their large raw split files were not included in the
source handoff, and the supplied review did not replay all that
arithmetic. They are retained as historical supplementary records, not
additional dependencies of the finite boundary theorem. Winner summaries
are not certificates of all terminal-pair isomorphism classes. Likewise,
a pendant-plus-core description applies only where the edge count
permits a pendant vertex, not to the entire dense range.

\subsection{A small certificate readers can
inspect}\label{a-small-certificate-readers-can-inspect}

At \(N=7\), \(\varepsilon=1\), the first excluded value is \(d=2\). This
is one of the five exceptions to the original sharpness assertion. With
\(q=1-p\), \[\Delta(p)=(1-q^3)\operatorname{Rel}(C_7(2);p)
-p\operatorname{Rel}(C_7(4);p)
=\sum_{j=0}^{20}D_jp^jq^{20-j}.\] Only \(D_7,\ldots,D_{19}\) are
nonzero, and they are
\[840,\ 7640,\ 19904,\ 30348,\ 32615,\ 26418,\ 16510,\]
\[7976,\ 2938,\ 800,\ 152,\ 18,\ 1.\] They are all positive, so
\(\Delta(p)>0\) throughout \(0<p<1\); at \(p=1/2\) it is \(9135/65536\).
A standalone enumerator checks all \(2^{17}+2^{19}=655360\) edge sets
without importing the submitted graph routines. This example also
explains why excluded parameter values must be tested explicitly rather
than inferred from a loop over the required positive range.

\subsection{Reproducibility boundary}\label{reproducibility-boundary}

The package's claim-to-file map gives the exact scripts, corpus files
and logs. Required checks return nonzero on a failed comparison, and
negative controls exercise the loop convention and incorrect boundary
assertions. Historical logs are not rewritten. Rebuild the C
implementation from source; bundled platform-specific executables are
not part of the public evidence contract. The integer bounds apply only
in the documented supported range. Reproducing outputs checks the
implemented finite statement, not the imported literature or every step
of the universal prose proof.

\section{Discussion}\label{discussion}

The threshold \(m=\binom n2-n\) has a transparent meaning in the light
of Lemma \ref{lem:identity}: a split event is a connectivity event on
the identified graph \(H/st\) minus a connectivity event on \(H\)
itself, and the balloon candidate is simultaneously the graph whose
identification is most reliable (Lemma \ref{lem:identified}, true for
all \(k\le n\)) and the graph that is itself least reliable
(Petingi--Saccoman--Schoppmann, known for \(k\le n-2\)). For
\(k\ge n+1\) the first property fails to compensate: the
lexicographically forced candidate is beaten at the \(2\)-forest
coefficient by graphs whose missing edges are spread rather than
concentrated, as Section 3 shows with explicit spanning-tree counts.

The two edge counts left conditional, \(m=\binom{n-1}2\) and
\(m=\binom{n-1}2-1\), are exactly the first two below the range in which
the all-terminal ``uniformly least reliable'' problem is solved. The
statement that would complete the characterisation is the following
instance of the conjecture of Boesch, Satyanarayana and Suffel.

\begin{conjecture}\label{conj:boundary}
Let $n\ge 7$ and let $H$ be a connected graph with $n$ vertices and $m$ edges.
\begin{itemize}
\item[(a)] If $m=\binom{n-1}2$, then $\mathrm{Rel}(H;p)\ge p\,\mathrm{Rel}(K_{n-1}-e;p)$ for all $p\in[0,1]$, where $K_{n-1}-e$ is the complete graph with one edge removed.
\item[(b)] If $m=\binom{n-1}2-1$, then $\mathrm{Rel}(H;p)\ge p\,\mathrm{Rel}(K_{n-1}-P_3;p)$ for all $p\in[0,1]$, where $K_{n-1}-P_3$ is the complete graph with two adjacent edges removed.
\end{itemize}
\end{conjecture}

The right-hand sides are the reliabilities of the balloon graphs
\(B_{n,m}\) at the two edge counts (a pendant vertex attached anywhere
to the core contributes the factor \(p\), so the choice of attachment
vertex is immaterial). By Theorem \ref{thm:Bprime}, Conjecture
\ref{conj:boundary} implies that a UMSR graph exists in
\(\mathcal T_{n,m}\) if and only if \(m=n-1\) or \(m\ge\binom n2-n\),
for every \(n\ge 7\). Theorem \ref{thm:reduction} shows that Conjecture
\ref{conj:boundary} for \(n\ge11\) follows from the inequality
\((\star)_{n-1,\varepsilon}\) between the two cone graphs
\(C_{n-1}(n-3-\varepsilon)\) and \(C_{n-1}(n-3-\varepsilon-d)\), and
Corollary \ref{cor:n41} establishes it for \(7\le n\le 51\). A
sufficient remaining research target is therefore a family indexed by
\(N,\varepsilon,d\), uniformly quantified over \(p\):

\begin{conjecture}\label{conj:star}
$(\star)_{N,\varepsilon}$ holds for every $N\ge 7$ and $\varepsilon\in\{0,1\}$; that is, for $h=N-2-\varepsilon$ and $1\le d\le N-5-\varepsilon$,
$$(1-q^{d+1})\,\mathrm{Rel}\bigl(K_N-K_{1,1+\varepsilon+d};p\bigr)\ \ge\ p\,\mathrm{Rel}\bigl(K_N-K_{1,1+\varepsilon};p\bigr)\qquad(0\le p\le 1).$$
\end{conjecture}

Conjecture \ref{conj:star} implies Conjecture \ref{conj:boundary} and
hence completes the characterisation for all \(n\ge7\). It holds
coefficient-wise for \(N\le50\) (Proposition \ref{prop:star}). Theorem
\ref{thm:coneintervals} proves it at every order for \(p\ge1/2\) and on
an explicit interval near zero; only the intervening interval remains to
be proved at unrestricted orders. The first excluded parameter value has
a leading-coefficient obstruction for \(N\ge11\), so a proof cannot
simply extend the required range without qualification.

The theorem classifies existence, not all terminal-preserving
isomorphism classes of optimisers. Romero's theorem identifies the
optimisers, whenever they exist, by split-equivalence to the balloon
polynomial. The supplied winner summaries retain only representative
terminal pairs and do not certify an enumeration of every
terminal-placement class. We make no new terminal-isomorphism uniqueness
claim on that evidence.

\section*{References}\label{references}
\addcontentsline{toc}{section}{References}

\textbf{Split reliability and uniformly optimal reliability (primary
sources)}

Brown, J. I., \& McMullin, I. (2023). On the split reliability of
graphs. \emph{Networks, 82}(2), 177--185.
https://doi.org/10.1002/net.22166 (preprint:
https://arxiv.org/abs/2306.03884)

Romero, P. (2025). Characterization of locally most split reliable
graphs. \emph{Theoretical Computer Science, 1047}, 115327.
https://doi.org/10.1016/j.tcs.2025.115327 (preprint:
https://arxiv.org/abs/2503.15453)

Romero, P. (2022). Uniformly optimally reliable graphs: A survey.
\emph{Networks, 80}(4), 466--481. https://doi.org/10.1002/net.22085

Petingi, L., Saccoman, J. T., \& Schoppmann, L. (1996). Uniformly least
reliable graphs. \emph{Networks, 27}(2), 125--131. (Abstract and
bibliographic record: https://iaorifors.com/paper/34598)

Romero, P. (2022). Universal reliability bounds for sparse networks.
\emph{IEEE Transactions on Reliability, 71}(1), 359--369. First
published online in 2021. https://doi.org/10.1109/TR.2021.3061075 (open
access copy:
https://redi.anii.org.uy/jspui/bitstream/20.500.12381/648/1/10\%20(1).pdf)

Brown, J. I., \& Mol, L. (2017). On the roots of all-terminal
reliability polynomials. \emph{Discrete Mathematics, 340}(6),
1287--1299. https://arxiv.org/abs/1703.10566

\textbf{Reliability-monotone graph transformations}

Brown, J. I., Colbourn, C. J., \& Devitt, J. S. (1993). Network
transformations and bounding network reliability. \emph{Networks,
23}(1), 1--17. https://doi.org/10.1002/net.3230230103

Satyanarayana, A., Schoppmann, L., \& Suffel, C. L. (1992). A
reliability-improving graph transformation with applications to network
reliability. \emph{Networks, 22}(2), 209--216.
https://doi.org/10.1002/net.3230220206

Kelmans, A. K. (1981). On graphs with randomly deleted edges. \emph{Acta
Mathematica Academiae Scientiarum Hungaricae, 37}(1--3), 77--88.
https://doi.org/10.1007/BF01904874 (statement used: Csikvári, 2011,
Theorem 2.2)

Csikvári, P. (2011). Applications of the Kelmans transformation:
Extremality of the threshold graphs. \emph{Electronic Journal of
Combinatorics, 18}(1), P182. https://doi.org/10.37236/669

\textbf{Software}

McKay, B. D., \& Piperno, A. (2014). Practical graph isomorphism, II.
\emph{Journal of Symbolic Computation, 60}, 94--112.
https://doi.org/10.1016/j.jsc.2013.09.003 (nauty, including
\texttt{geng}: https://pallini.di.uniroma1.it/)

\end{document}
